\chapter{Introduction}
\label{chap:introduction}

Satellite communication links are affected by propagation impairments that can
temporarily reduce the received signal quality. This problem is especially
relevant at high frequencies, where atmospheric phenomena such as rain can
produce strong fades on Earth-space paths \cite{panagopoulos2004satellite,
iturP618}. In operational systems, a fade is not only a physical phenomenon to
be described after it occurs. It is also a decision problem: when the primary
link degrades, an operator or automated controller may need to activate a
backup communication resource. Activating this resource too late may lead to
service degradation, while activating it too often or for too long may waste
capacity and increase operational cost.

The central question of this thesis is therefore not simply whether the future
satellite signal can be predicted. The more relevant question is whether
information available at the current time can support a timely and reliable
switch decision. This distinction changes how the problem should be modelled
and evaluated. A model with a low forecast error is not necessarily a good
switch controller if its errors occur around the operational threshold. A
model that estimates a duration or a probability may be more useful than a
trajectory forecaster if that quantity is more directly connected to the
decision. For this reason, the thesis treats satellite fade nowcasting as a
decision-oriented time-series problem.

\section{Operational Motivation}
\label{sec:introduction-operational-motivation}

Fade mitigation techniques are commonly used to preserve link availability
under adverse propagation conditions. Depending on the system, mitigation may
involve power control, adaptive coding and modulation, site diversity, or the
activation of an alternative link. The specific mitigation mechanism is outside
the scope of this thesis. The focus here is the upstream decision that such a
mechanism requires: given the signal observed up to the current timestamp,
should the system switch to a backup resource?

This decision has an intrinsically temporal structure. A very short threshold
crossing may not justify a switch, because the system could recover before the
backup activation becomes useful. Conversely, a fade that persists for several
minutes may justify switching even if the instantaneous signal is not at its
worst point yet. The operational objective is therefore tied to persistence,
residual duration, and event evolution, not only to the instantaneous signal
level. This is consistent with the broader propagation literature, where fade
dynamics and time-series behavior are treated as relevant design quantities
for Earth-space communication systems \cite{iturP1623,iturP1853}.

This thesis was carried out at M.B.I. S.r.l. \cite{mbiGroup} within the broader
context of the NEFOCAST project, for which M.B.I. is the lead company. NEFOCAST
investigates an integrated technological platform for estimating rainfall
fields from satellite-signal attenuation measurements
\cite{nefocastProject}. This industrial and research setting provides the
application context for the
decision-oriented nowcasting problem addressed in this work.

The data used in this work consist of signal time series acquired at the
Fucino ground station from multiple files and time periods. The available files
are heterogeneous: some are uplink or downlink beacon measurements, some are
associated with different satellites, and only a subset contains auxiliary
columns. To make the experimental comparison reproducible and fair, the thesis
adopts a strict signal-only policy. Every dataset is standardized to a common
time and signal representation, and all models use only the historical signal
as input. This avoids comparing models that have access to different physical
variables and keeps the focus on what can be learned from the signal evolution
itself.

\section{Satellite Link Fade Nowcasting}
\label{sec:introduction-nowcasting-problem}

Nowcasting refers to short-term prediction at or near the present time. In this
work, the prediction time is a timestamp \(t\) in a satellite signal time
series. The information available to a model is the recent past and present
signal history, while the target concerns the future behavior of the same
signal or of the ongoing fade event. No future observation can be used as an
input. This no-leakage constraint is essential because the final goal is an
operational decision that could be made online.

The thesis investigates four task formulations that all serve the same final
switching objective but estimate different intermediate quantities:

\begin{itemize}
  \item autoregressive forecasting, where the model predicts a future signal
  trajectory;
  \item current-level persistence, where the model estimates how long the
  signal will remain close to its current level;
  \item long-fade detection, where the model estimates whether the current
  grouped fade event will last longer than an operational duration threshold;
  \item survival persistence, where the model estimates a residual-duration
  survival function for the ongoing fade event.
\end{itemize}

These formulations are not interchangeable. They answer different operational
questions and produce different native outputs: trajectories, scalar
durations, class probabilities, or survival curves. To compare them, the
outputs must be converted into binary switch sequences and evaluated against a
common reference. This thesis therefore builds a shared Perfect Switch
reference, applies common switch post-processing, and reports operational
metrics in addition to native model metrics.

\section{Research Questions}
\label{sec:introduction-research-questions}

The work is guided by four research questions.

\begin{enumerate}
  \item \textbf{Can task-specific machine-learning models improve operational
  switch decisions for satellite link fades?} The objective is to evaluate
  whether learned models produce switch sequences that are closer to the
  reference decision than simple threshold reasoning.

  \item \textbf{Which predictive formulation is most useful for the switching
  objective?} The thesis compares trajectory forecasting, duration regression,
  event-level classification, and survival modeling because the best
  formulation is not obvious a priori.

  \item \textbf{Do native predictive metrics explain operational usefulness?}
  Forecast RMSE, duration MAE, classification AUPRC, and survival metrics are
  informative, but the final decision depends on thresholding,
  post-processing, and event timing. The thesis therefore studies how native
  metrics relate to switch metrics.

  \item \textbf{How important are common reference construction and
  post-processing?} Since different tasks produce outputs on different
  timestamp domains and with different semantics, the comparison must separate
  modeling differences from evaluation artifacts.
\end{enumerate}

These questions deliberately place the switch decision at the centre of the
study. The intermediate predictions remain important, but they are evaluated
as components of a decision pipeline rather than as isolated forecasting
outputs.

\section{Contributions}
\label{sec:introduction-contributions}

The main methodological and empirical contributions of the thesis are the
following.

\begin{enumerate}
  \item It formulates satellite fade nowcasting as a decision-oriented problem
  through four complementary predictive tasks: autoregressive forecasting,
  current-level persistence, long-fade detection, and survival persistence.
  Each task estimates a different intermediate quantity but is connected
  explicitly to the same operational switch objective.

  \item It establishes a shared switch-evaluation framework with a common
  Perfect Switch reference, shared post-processing, task-native comparisons,
  and a cross-task reference-grid comparison. This framework makes models with
  different native outputs operationally comparable.

  \item It evaluates a range of model families under these formulations,
  including GRU models, PatchTST, Chronos zero-shot forecasting, XGBoost
  baselines, learnable shapelet models, temporal convolutional networks,
  XGBoost-AFT, and a discrete-time survival TCN.

  \item It provides an empirical analysis of the relationship between native
  prediction quality and operational switch quality. The results show that the
  best native predictor is not always the best switch controller and that task
  formulation strongly affects the precision-recall trade-off.

  \item To support reproducibility, it develops a signal-only preprocessing
  pipeline for heterogeneous satellite fade and beacon files. The pipeline
  validates and standardizes raw data, handles gaps conservatively, constructs
  task-specific datasets, and preserves the metadata required for event-level
  evaluation.
\end{enumerate}

The source-code repository for the experimental framework is hosted on GitHub:
\url{https://github.com/Giusgarus/Nowcasting-project}.
It contains the task implementations, experiment configurations, and evaluation
scripts used in this work.

\section{Thesis Structure}
\label{sec:introduction-structure}

\Cref{chap:background} introduces the technical background needed for the
work. It covers time-series data, machine-learning models used in the thesis,
survival analysis, nowcasting and validation principles, and the satellite
communication context.

\Cref{chap:data-reference-framework} describes the data preparation and
reference-decision framework. It explains the available signal files, the
signal-only policy, parsing and validation, sampling gaps, conservative
imputation, event construction, the external holdout split, the Perfect Switch
reference, and the shared switch post-processing rules.

\Cref{chap:task-formulations} presents the four modeling formulations. For
each formulation, it defines the input representation, target construction,
model families, native output, and switch-conversion rule.

\Cref{chap:experimental-setup} describes the experimental protocol, including
dataset construction, train-validation-test usage, hyperparameter search,
native metrics, switch metrics, intra-task comparisons, cross-task
comparisons, and reproducibility settings.

\Cref{chap:results} reports and discusses the empirical results. It first
analyses the native predictive performance within each task and then focuses
on the switch metrics and event-level behavior that determine operational
usefulness.

\Cref{chap:conclusions} summarizes the main findings, answers the research
questions, discusses limitations, and outlines future work.
